% Prose rendering of PrimeTensor/Analysis/Control.lean.
% Fragment: no \documentclass, no preamble, no document environment.

\section{Quantitative scale control for tangent maps}
\label{Analysis:Control:sec}

\leanfile{PrimeTensor/Analysis/Control.lean}

\subsection*{What this file establishes}

Every tangent map so far has been an arbitrary morphism of the carrier.
\texttt{PrimeTensor/Analysis/Derivative.lean} observed that nothing required one
to be continuous, and that nothing up to that point needed it.  The chain rule
does, and the module header says why the obvious repair is insufficient:
ordinary continuity at the pivot is not enough, because the derivative measures
a \emph{relative} rate, so what is wanted is a quantitative statement --- an
input refinement gain that delivers a prescribed output refinement gain,
uniformly in the ambient level.  That property is $\name{ScaleControlled}$, and
this file introduces it, verifies it for every tangent map the development has
built, and proves the two consequences the chain rule will consume.

Getting there requires a fact outstanding since
\texttt{PrimeTensor/Scale.lean}: that the levels are \emph{nested}, so that
closeness at a fine level implies closeness at a coarse one.  That file
recorded the statement as true and unproved; this one proves it, and
Remark~\ref{Analysis:Control:rem:why-now} explains why it becomes unavoidable
exactly here.

\subsection{The levels are nested}

\begin{lemma}[\lean{MulRat.scaleWithin\_succ\_weaken}]
\label{Analysis:Control:lem:weaken-rat}
$\name{ScaleWithin}(\ctor{succ}\,\ell, a, b)$ implies
$\name{ScaleWithin}(\ell, a, b)$.
\end{lemma}

\begin{leancode}
theorem scaleWithin_succ_weaken {level : Depth} {a b : MulRat}
    (h : ScaleWithin (.succ level) a b) :
    ScaleWithin level a b := by
  unfold ScaleWithin at h ⊢

  have hone :
      (1 : MulRat) * (1 : MulRat) < two := by
    rw [one_mul]
    exact one_lt_two
\end{leancode}

\begin{proof}
Write $q = \name{scalePow}(\name{ratio}(a,b), \ell)$, so that the forward
hypothesis is $q \cdot q < \name{two}$.  Also $1 \cdot 1 < \name{two}$, by
cancelling a unit and citing $\name{one\_lt\_two}$.  Now apply
$\name{mul\_lt\_two\_of\_sq\_lt\_two}$ of
\texttt{PrimeTensor/Scale/Laws.lean} to the pair $(q, 1)$: both have square
below the ruler, so their product $q \cdot 1$ does, and cancelling the unit
gives $q < \name{two}$, which is the forward half of the conclusion.  The
backward half is the same argument on the reversed ratio.
\end{proof}

\begin{remark}[Monotonicity, and the route taken to it]
\label{Analysis:Control:rem:monotone}
\texttt{PrimeTensor/Scale.lean} stated this as something the ladder ought to
satisfy, gave the informal argument --- if $r^{2^{k-1}} \geq 2$ then
$r^{2^{k}} \geq 4 > 2$ --- and noted that the argument needed a comparison fact
not then available.  \texttt{PrimeTensor/Scale/Laws.lean} supplied that fact
and observed the statement had become a corollary of the composition law with
the middle magnitude taken to be $a$ itself.

The proof above takes neither route.  It pairs the hypothesis not with a second
hypothesis but with the trivial bound $1 \cdot 1 < \name{two}$, and lets the
arithmetic core do the rest.  The effect is the same and the shape is more
direct: weakening a level is composing with the identity comparison, so it
needs no new inequality at all --- only the observation that the pivot is
itself within the ruler, which was the very first lemma of
\texttt{PrimeTensor/Scale.lean}.
\end{remark}

\begin{lemma}[\lean{MulReal.scaleNear\_succ\_weaken} and \lean{scaleNear\_weaken}]
\label{Analysis:Control:lem:weaken-real}
$\name{ScaleNear}(\ctor{succ}\,\ell, a, b)$ implies
$\name{ScaleNear}(\ell, a, b)$; and more generally, if
$\name{AtOrAfter}(c, f)$ then $\name{ScaleNear}(f,a,b)$ implies
$\name{ScaleNear}(c,a,b)$.
\end{lemma}

\begin{leancode}
theorem scaleNear_weaken {coarse fine : Depth} {a b : MulReal}
    (hcf : Depth.AtOrAfter coarse fine)
    (h : ScaleNear fine a b) :
    ScaleNear coarse a b := by
  induction hcf with
  | here =>
      exact h
  | later hcf ih =>
      exact ih (scaleNear_succ_weaken h)
\end{leancode}

\begin{proof}
The first transports Lemma~\ref{Analysis:Control:lem:weaken-rat} across the
quotient, keeping the same witnessing streams and the same anchor and applying
it termwise.  The second iterates: induction on the proof that the fine level
lies at or after the coarse one peels off one $\ctor{succ}$ at a time, each
peel being one application of the first.
\end{proof}

So the family $\{\name{ScaleNear}(\ell,-,-)\}$ is decreasing along
$\name{AtOrAfter}$ --- the levels really are a nested system of tolerances, and
the order on $\name{Depth}$ built in
\texttt{PrimeTensor/Completion.lean} for indexing tails now also orders
tolerances.

\subsection{Advancing respects the order}

\begin{lemma}[\lean{Depth.advance\_atOrAfter}, \lean{advance\_left\_join} and
  \lean{advance\_right\_join}]\label{Analysis:Control:lem:advance}
If $\name{AtOrAfter}(a,b)$ then
$\name{AtOrAfter}(\name{advance}(\ell,a), \name{advance}(\ell,b))$.
Consequently $\name{advance}(\ell, \name{join}(a,b))$ lies at or after both
$\name{advance}(\ell,a)$ and $\name{advance}(\ell,b)$.
\end{lemma}

\begin{proof}
The first is induction on the given proof, each $\ctor{later}$ becoming one
$\ctor{later}$ on the advanced pair, since $\name{advance}$ recurses on its
second argument by a single $\ctor{succ}$.  The two corollaries apply it to the
bounds on $\name{join}$ from \texttt{PrimeTensor/Completion.lean}.
\end{proof}

\begin{remark}[Why monotonicity becomes unavoidable in this file]
\label{Analysis:Control:rem:why-now}
Six earlier proofs have combined two hypotheses by taking a $\name{join}$, and
none of them needed the levels to be nested.  The reason is that in all six the
join was taken in the \emph{anchor} coordinate: an index lying in the tail at
$\name{join}(c_1,c_2)$ lies in the tail at $c_1$ by transitivity of
$\name{AtOrAfter}$, and transitivity is all that is required.

Theorem~\ref{Analysis:Control:thm:mul} takes a join in the \emph{level}
coordinate instead --- of two source gains.  The hypothesis then arrives at
$\name{advance}(\ell, \name{join}(s_1,s_2))$, which is \emph{finer} than the
$\name{advance}(\ell, s_1)$ that the first control hypothesis expects, and a
finer nearness is not automatically a coarser one.  Converting it is exactly
Lemma~\ref{Analysis:Control:lem:weaken-real}, applied through
Lemma~\ref{Analysis:Control:lem:advance}.  That is why the statement
outstanding since \texttt{PrimeTensor/Scale.lean} is proved here and not
earlier: this is the first place a join happens among tolerances rather than
among stages.
\end{remark}

\subsection{Scale control}

\begin{definition}[\lean{MulTangentMap.ScaleControlled}]
\label{Analysis:Control:def:controlled}
A tangent map $D$ is \emph{scale controlled} when
\[
  \forall\, t\;
  \exists\, s\;
  \forall\, \ell\; \forall\, h : \quad
  \name{ScaleNear}(\name{advance}(\ell,s),\, h,\, 1)
  \;\longrightarrow\;
  \name{ScaleNear}(\name{advance}(\ell,t),\, D(h),\, 1).
\]
\end{definition}

\begin{leancode}
def ScaleControlled (D : MulTangentMap) : Prop :=
  ∀ targetGain : Depth,
    ∃ sourceGain : Depth,
      ∀ level : Depth, ∀ h : MulReal,
        MulReal.ScaleNear (Depth.advance level sourceGain) h 1 →
        MulReal.ScaleNear (Depth.advance level targetGain) (D h) 1
\end{leancode}

\begin{designnote}[What the quantifier order buys]
\label{Analysis:Control:note:uniform}
The source gain $s$ is chosen \emph{before} the ambient level $\ell$, so one
$s$ must work at every level simultaneously.  Had the order been
$\forall t\,\forall \ell\,\exists s$, the modulus could depend on the ambient
scale and the definition would be a level-by-level continuity statement.  As
written it is uniform, and that uniformity is what the docstring calls finite
scale distortion.

In the logarithmic chart of
\texttt{PrimeTensor/Analysis/Derivative.lean} the reading is the familiar one.
Writing $T = 2^{-(\size{\ell}-1)}$ for the tolerance at level $\ell$, the
hypothesis is $\lvert \log h \rvert < 2^{-\size{s}}\,T$ and the conclusion is
$\lvert \log D(h) \rvert < 2^{-\size{t}}\,T$, for every dyadic $T$.  Since the
induced map on logarithms is additive, that says it is bounded, with
$2^{\size{s}-\size{t}}$ serving as a bound on its operator norm.  So
$\name{ScaleControlled}$ is boundedness of the tangent map --- exactly the
hypothesis \texttt{PrimeTensor/Analysis/Derivative.lean} declined to impose,
imposed now that something needs it.
\end{designnote}

\begin{lemma}[\lean{scaleControlled\_identity} and \lean{scaleControlled\_trivial}]
\label{Analysis:Control:lem:examples}
$\name{identity}$ and $\name{trivial}$ are scale controlled.
\end{lemma}

\begin{proof}
For the identity take $s = t$ and return the hypothesis unchanged --- no
distortion at all.  For the trivial map take $s = t$ as well; the value is the
pivot, and $\name{scaleNear\_refl}$ places it near itself at every level, so
the hypothesis is discarded.
\end{proof}

\begin{theorem}[\lean{ScaleControlled.inv}]\label{Analysis:Control:thm:inv}
If $D$ is scale controlled then so is $D^{-1}$, with the same source gain.
\end{theorem}

\begin{proof}
Take the source gain $D$ supplies for the same target, apply $D$'s control, and
invert the conclusion with $\name{scaleNear\_inv}$ of
\texttt{PrimeTensor/Analysis/Near/Laws.lean}; the pivot is fixed by inversion,
so the right-hand side stays $1$.
\end{proof}

\begin{theorem}[\lean{ScaleControlled.mul}]\label{Analysis:Control:thm:mul}
If $D$ and $E$ are scale controlled then so is $D \cdot E$.
\end{theorem}

\begin{leancode}
theorem ScaleControlled.mul {D E : MulTangentMap}
    (hD : ScaleControlled D)
    (hE : ScaleControlled E) :
    ScaleControlled (mul D E) := by
  intro targetGain

  obtain ⟨sourceD, hcontrolD⟩ := hD (.succ targetGain)
  obtain ⟨sourceE, hcontrolE⟩ := hE (.succ targetGain)

  let sourceGain := Depth.join sourceD sourceE
\end{leancode}

\begin{proof}
Given a target gain $t$, ask each factor for $\ctor{succ}\,t$ --- one deeper,
because the product will spend a level --- obtaining source gains $s_D$ and
$s_E$, and offer $\name{join}(s_D, s_E)$ as the source gain.

Let $h$ be near the pivot at $\name{advance}(\ell, \name{join}(s_D,s_E))$.  By
Lemma~\ref{Analysis:Control:lem:advance} that level is at or after both
$\name{advance}(\ell,s_D)$ and $\name{advance}(\ell,s_E)$, so
Lemma~\ref{Analysis:Control:lem:weaken-real} weakens the hypothesis to each of
them and both controls apply, putting $D(h)$ and $E(h)$ near the pivot at
$\name{advance}(\ell, \ctor{succ}\,t)$.  That level is
$\ctor{succ}\,\name{advance}(\ell,t)$ by $\name{advance\_succ}$, so
$\name{scaleNear\_mul}$ combines the two at $\name{advance}(\ell,t)$; the
right-hand side is $1 \cdot 1$, which cancels.
\end{proof}

\begin{corollary}[\lean{scaleControlled\_depthPow}]
\label{Analysis:Control:cor:pow}
Every positive-depth tangent power $\name{depthPow}(d)$ is scale controlled.
\end{corollary}

\begin{proof}
Structural recursion: at $\ctor{one}$ it is the identity, and at
$\ctor{succ}\,d$ it is the product of the identity with the recursive case.
\end{proof}

\begin{remark}[Controlled is a property of the family, not of tangent maps]
\label{Analysis:Control:rem:not-all}
No theorem here says that a tangent map is scale controlled; the property is
verified one construction at a time.  It could not be otherwise --- a tangent
map is any morphism of the carrier whatever, and
\texttt{PrimeTensor/Analysis/Derivative.lean} deliberately imposed no
regularity, so a wild one is not excluded by anything.  What
Lemma~\ref{Analysis:Control:lem:examples},
Theorems~\ref{Analysis:Control:thm:inv} and~\ref{Analysis:Control:thm:mul} and
Corollary~\ref{Analysis:Control:cor:pow} establish jointly is that the
\emph{closure} of $\{\name{identity}, \name{trivial}\}$ under product and
inverse consists of controlled maps --- which, by
\texttt{PrimeTensor/Analysis/Examples.lean}, is every tangent map the
development has so far produced as a derivative.  The tangent ratio is not
separately checked, although it is a product with an inverse and so is covered.
\end{remark}

\subsection{The two consequences}

\begin{theorem}[\lean{ScaleControlled.maps\_converges}]
\label{Analysis:Control:thm:converges}
A scale-controlled $D$ sends pivot-convergent sequences to pivot-convergent
sequences.
\end{theorem}

\begin{proof}
Fix a level $\ell$.  Ask the control for the smallest target gain,
$\ctor{one}$, obtaining a source gain $s$, and ask the convergence hypothesis
for the level $\name{advance}(\ell, s)$, obtaining an anchor.  Beyond it, the
control puts $D(s_n)$ near the pivot at $\name{advance}(\ell, \ctor{one})$,
which is $\ctor{succ}\,\ell$; weakening once by
Lemma~\ref{Analysis:Control:lem:weaken-real} gives level $\ell$, as required.
\end{proof}

This is continuity at the pivot, in the sequential form of
\texttt{PrimeTensor/Analysis/Limit.lean}.  The implication runs one way only:
control is defined uniformly in the level, and nothing here derives it back
from the sequential statement.  That gap is the content of the header's remark
that ordinary continuity is not enough.

\begin{theorem}[\lean{negligible\_map}]\label{Analysis:Control:thm:negligible-map}
If $D$ is scale controlled and $\name{error}$ is negligible relative to
$\name{base}$, then $n \mapsto D(\name{error}_n)$ is negligible relative to
$\name{base}$.
\end{theorem}

\begin{leancode}
theorem negligible_map
    {D : MulTangentMap}
    (hD : D.ScaleControlled)
    {error base : MulReal.Seq}
    (herror : NegligibleRelative error base) :
    NegligibleRelative (fun n => D (error n)) base := by
  intro targetGain

  obtain ⟨sourceGain, hcontrol⟩ := hD targetGain
  obtain ⟨anchor, htail⟩ := herror sourceGain

  refine ⟨anchor, ?_⟩
  intro n hn level hbase
\end{leancode}

\begin{proof}
Given a target gain, take the source gain the control supplies and the anchor
the negligibility supplies at that source gain.  Beyond the anchor, whatever
level the base achieves, the error is near the pivot at the source-advanced
level, and the control converts that into the target-advanced level for
$D(\name{error}_n)$.  The two definitions match term for term; the proof is
their composition and nothing more.
\end{proof}

\begin{remark}[Everything here is for the chain rule]
\label{Analysis:Control:rem:purpose}
Theorem~\ref{Analysis:Control:thm:negligible-map} is the reason the file
exists.  A chain rule must push the error of the inner function through the
outer tangent map and still call it negligible, and negligibility is a
statement about levels, so the outer map has to move levels predictably.
Nothing in this file states a chain rule, composes two functions, or defines
composition of tangent maps; all of that is downstream.  What is settled here
is the hypothesis such a rule will carry.
\end{remark}

\begin{remark}[What is not proved]\label{Analysis:Control:rem:missing}
$\name{ScaleControlled}$ is not shown closed under composition of tangent maps,
composition itself not being defined.  There is no converse to
Theorem~\ref{Analysis:Control:thm:converges}, and no example of a tangent map
that fails to be controlled --- so the notion is not shown to be a genuine
restriction, only a sufficient one.  Nothing relates control to
$\name{HasMulDerivativeAt}$: a derivative is not required to be controlled, and
a function is not shown to have a controlled derivative.
\end{remark}

\subsection*{Summary of the correspondence}

\begin{center}
\small
\begin{tabular}{@{}lll@{}}
\hline
Lean declaration & Kind & Here \\
\hline
\lean{MulRat.scaleWithin\_succ\_weaken}   & theorem    & Lem.~\ref{Analysis:Control:lem:weaken-rat} \\
\lean{MulReal.scaleNear\_succ\_weaken}    & theorem    & Lem.~\ref{Analysis:Control:lem:weaken-real} \\
\lean{MulReal.scaleNear\_weaken}          & theorem    & Lem.~\ref{Analysis:Control:lem:weaken-real} \\
\lean{Depth.advance\_atOrAfter}           & theorem    & Lem.~\ref{Analysis:Control:lem:advance} \\
\lean{Depth.advance\_left\_join}          & theorem    & Lem.~\ref{Analysis:Control:lem:advance} \\
\lean{Depth.advance\_right\_join}         & theorem    & Lem.~\ref{Analysis:Control:lem:advance} \\
\lean{MulTangentMap.ScaleControlled}      & definition & Def.~\ref{Analysis:Control:def:controlled} \\
\lean{scaleControlled\_identity}          & theorem    & Lem.~\ref{Analysis:Control:lem:examples} \\
\lean{scaleControlled\_trivial}           & theorem    & Lem.~\ref{Analysis:Control:lem:examples} \\
\lean{ScaleControlled.inv}                & theorem    & Thm.~\ref{Analysis:Control:thm:inv} \\
\lean{ScaleControlled.mul}                & theorem    & Thm.~\ref{Analysis:Control:thm:mul} \\
\lean{scaleControlled\_depthPow}          & theorem    & Cor.~\ref{Analysis:Control:cor:pow} \\
\lean{ScaleControlled.maps\_converges}    & theorem    & Thm.~\ref{Analysis:Control:thm:converges} \\
\lean{MulDifferential.negligible\_map}    & theorem    & Thm.~\ref{Analysis:Control:thm:negligible-map} \\
\hline
\end{tabular}
\end{center}

\noindent
Every declaration of \texttt{PrimeTensor/Analysis/Control.lean} appears above.
The file contains no \lean{sorry}, no axiom, and no unproved named proposition.
Design note~\ref{Analysis:Control:note:uniform} and
Remarks~\ref{Analysis:Control:rem:monotone},
\ref{Analysis:Control:rem:why-now}, \ref{Analysis:Control:rem:not-all},
\ref{Analysis:Control:rem:purpose} and~\ref{Analysis:Control:rem:missing} are
stated for the reader and are not declarations of the file; the boundedness
reading of scale control is a gloss and is not formalised.
